% Imporditud: tools/import_eksperimendid.py
% Allikas: Eesti füüsikaolümpiaadi eksperimendi ülesannete
%   kogu (Rael Kalda, 2023), ülesanne Ü33, lahendus L33.
\setAuthor{EFO žürii}
\setRound{piirkonnavoor}
\setYear{1996}
\setNumber{K E2}
\setDifficulty{2}
\setTopic{termodünaamika}
\setVahendid{termomeeter, kalorimeeter veega, teadaoleva massi ja erisoojusega keha temperatuuriga $\SI{100}{\celsius}$}
\setPunktid{10}
\setAllikas{Eksperimendi kogu Ü33, lahendus lk 47}
\setLahendusseis{toores}

\prob{Vee mass}
% Ainus ülesanne kogumikus, millel EI OLE eraldi rida "Katsevahendid:".
% Vahendid on loetletud ülesande tekstis endas ja ülal olev loend on sealt
% välja kirjutatud, midagi juurde lisamata.
Kasutades termomeetrit, määrata kalorimeetris oleva vee mass. Abivahendina
saab kasutada tuntud massiga keha, mille temperatuur on
$\SI{100}{\celsius}$ ja mille erisoojus on teada. Samuti on teada
kalorimeetri mass ja erisoojus ning vee erisoojus. Hinnata mõõteviga.

\hint

\solu
Soojusliku tasakaalu võrrand:

cV mV (t$^\circ -$ t$^\circ$0) + cK mK (t$^\circ -$ t$^\circ$0) + cM cM ($\theta -$ t$^\circ$) = 0

Siit avaldame: mK = $-$ cM mM ($\theta -$ t$^\circ$) $-$ cV mV (t$^\circ -$ $t_0$$^\circ$) 2p cK mK (t$^\circ -$ $t_0$$^\circ$)

Kus t$^\circ$0 on vee ja kalorimeetri algtemperatuur, t$^\circ$ on vee ja kalorimeetri lõpptemperatuur, $\theta$ = 100$^\circ$C. Vea hindamine ja arvutamine teoreetiliselt (3 p). Kui praktiliselt eksib mitte üle 10\%, siis 5 p, kui üle 10\%, siis 3 p. Kui üle 20\%, siis 1 p.
\probend
